% APENDICE A
\section{Estrutura do Repositório de Conhecimento}

\begin{tcolorbox}[enhanced,breakable,colback=black!2,colframe=black!40,boxrule=0.5pt,arc=3pt,left=6pt,right=6pt,top=6pt,bottom=6pt]
\begin{Verbatim}[fontsize=\small,baselinestretch=1.1]
knowledge-base/
├── README.md
├── guidelines/
│   ├── template-solution
│   └── rules-of-the-base
├── solutions/
│   ├── logic/
│   ├── discrete-math/
│   ├── algebra/
│   ├── probability/
│   └── algorithm/
├── problems/
│   ├── logic/
│   ├── discrete-math/
│   ├── algebra/
│   ├── probability/
│   └── algorithm/
├── techniques/
├── explanations/
├── core-knowledge/
└── resources/
    ├── cheat-sheets/
    └── study-materials/
\end{Verbatim}
\end{tcolorbox}

\begin{description}
  \item[\texttt{knowledge-base/}]
    Raiz do repositório. Ponto de entrada para todo o conteúdo estruturado de aprendizado.

  \item[\texttt{README.md}]
    Visão geral do projeto, instruções de uso e guia de contribuição.

  \item[\texttt{guidelines/}]
    Padrões e regras que definem como o conteúdo deve ser escrito e organizado no
    repositório.
    \begin{itemize}
      \item \texttt{template-solution} --- Modelo padrão para formatar cada entrada de
        solução.
      \item \texttt{rules-of-the-base} --- Convenções de nomenclatura, critérios de
        qualidade e padrões editoriais.
    \end{itemize}

  \item[\texttt{solutions/}]
    Soluções detalhadas passo a passo para os problemas, organizadas por área matemática
    (\texttt{logic/}, \texttt{discrete-math/}, \texttt{algebra/}, \texttt{probability/},
    \texttt{algorithm/}).

  \item[\texttt{problems/}]
    Problemas para prática, ainda sem solução ou propostos como desafio. Espelha exatamente
    a estrutura de subpastas de \texttt{solutions/}.

  \item[\texttt{techniques/}]
    Estratégias reutilizáveis de resolução de problemas e métodos matemáticos (ex.:
    indução, substituição, prova por contradição).

  \item[\texttt{explanations/}]
    Textos conceituais que explicam o \emph{porquê} de teoremas, demonstrações e ideias,
    indo além do procedimento mecânico.

  \item[\texttt{core-knowledge/}]
    Definições fundamentais, axiomas e teoremas que servem de base para o restante
    do repositório.

  \item[\texttt{resources/}]
    Materiais de apoio e documentos de referência rápida.
    \begin{itemize}
      \item \texttt{cheat-sheets/} --- Folhas de fórmulas e resumos de regras para
        consulta rápida.
      \item \texttt{study-materials/} --- Leituras complementares, links externos e textos
        de referência curados.
    \end{itemize}
\end{description}
